\chapter{Discussione dei risultati}
\label{chap8}
In questo capitolo verranno discussi i risultati presentati nel capitolo precedente e come questi hanno definito le scelte per l'avanzamento del progetto di tesi.

\paragraph{Scelta dell'architettura principale.} Tra le tre architetture mostrate (MLP, CNN, GCN) è stata scelta la migliore, per quanto riguarda le performance ottenute sullo split di test, per poter proseguire con ulteriori test.
Come possiamo notare in Tabella~\ref{tab:model_confront}, la \emph{CNN} presenta delle performance migliori rispetto alle altre due architetture prosposte; inoltre la scelta è stata determinata, oltre che dalle performance, dalla grandezza della rete (riferito al peso in byte del file contenente i pesi della rete), e sotto questo punto di vista la \emph{CNN} prevale ancora.

\paragraph{Dataset \emph{original} vs \emph{ridotto}.} Con la seconda tipologia di test, si è cercato di verificare se il dataset di training ridotto (Sezione~\ref{pulizia-dataset}) portasse degli effettivi miglioramente sulle prestazioni del modello.
Come si può notare in Tabella~\ref{tab:orig_vs_reduced}, il test è stato inconcludente in quanto non si hanno miglioramenti nelle performance, anzi sono calate leggermente in alcuni casi. 
Questo test presenta una problematica che merita di essere evidenziata: il dataset di training è stato ripulito con lo stesso modello che stiamo verificando, questo significa che anche se abbiamo utilizzato il test esposto nella Sezione~\ref{pulizia-dataset}, non possiamo essere certi che tutte le registrazioni che sono state eliminate presentassero effettivamente un malposizionamento. L'unico modo per essere certi è tramite revisione umana, per mano di cardiologi esperti, di tutto il dataset. Solo in questo modo possiamo essere sicuri di avere un dataset pulito.

\paragraph{Dataset esterno}
Con questo test è stata verificata la capacità dei modelli di generalizzare su registrazioni provenienti da un altro dataset (Sezione~\ref{china_dataset}). In Tabella~\ref{tab:china_confronto} possiamo notare come vengano riconfermate le discussioni del paragrafo precedente, dove i modelli allenati con il dataset originale performano generalmente meglio, tranne per alcune classi (le stesse in cui i modelli \emph{reduced} avevano delle performance leggeremente migliori).

\paragraph{Dataset esterno modificato}
Questo è il test principale con il quale è possibile misurare le performance effettive del modello.
I risultati riportati in Tabella \ref{tab:china_mod_classi} mostrano un'accuratezza complessiva molto elevata in tutte le configurazioni (compresa tra il 97,98\% e il 99,96\%), un valore tuttavia poco informativo data la forte sproporzione tra le classi, come evidenziato dai volumi di \emph{TN} in Tabella \ref{tab:china_mod_confusione} (dell'ordine delle decine di migliaia contro poche centinaia di TP). È quindi sulla precisione, sul recall e sull'F1 della classe di scambio che vanno lette le reali capacità del modello.

Le configurazioni LA\_LL, RA\_LA, RA\_LL e V1\_V2 confermano un comportamento solido: RA\_LL ottiene il miglior risultato in assoluto (F1 $\sim$.978, con recall pari a 1.0 e un numero di falsi positivi molto contenuto, 19 in tutti i seed), mentre RA\_LA raggiunge anch'essa recall perfetto ma con una precisione leggermente inferiore (.93-.94), dovuta a un numero di falsi positivi lievemente più alto (29-30). LA\_LL si distingue invece per l'andamento opposto: precisione molto alta (.978) ma recall più contenuto (.93-.94), segno che il modello tende in questo caso a essere predire la classe di scambio solo quando è molto "sicuro", a scapito di qualche falso negativo (26-31).

Le configurazioni che coinvolgono le classi V2-V6 mostrano invece una forte diminuzione della precisione, pur mantenendo un recall quasi sempre pari o vicino a 1.0. Il caso più critico è V2\_V3, con una precisione di appena .33 (F1 $\sim$.496) a fronte di oltre 850 falsi positivi per seed (Tabella \ref{tab:china_mod_confusione}), a indicare una forte tendenza del modello a confondere segmenti di classi adiacenti con la classe di scambio. Le configurazioni V3\_V4, V4\_V5 e V5\_V6 seguono un andamento simile ma meno marcato, con precisione crescente al diminuire del numero di falsi positivi (rispettivamente $\sim$ .666, $\sim$.469 e $\sim$.38), suggerendo che l'errore del modello sia sistematico e correlato alla vicinanza temporale/posizionale tra le classi coinvolte, più che a un problema isolato di una singola configurazione.

L'introduzione della soglia di confidenza al 90\% (Tabella \ref{tab:china_mod_treshold}) conferma questa ipotesi, infatti, per tutte le configurazioni critiche si osserva un netto calo dei falsi positivi, con un miglioramento della precisione (V3\_V4 passa da $\sim$.666 a $\sim$.833, V4\_V5 da $\sim$.469 a $\sim$.617, V5\_V6 da $\sim$.38 a $\sim$.66), a fronte di una perdita minima o nulla di recall (al più 3 falsi negativi aggiuntivi, come in V5\_V6). Fa eccezione V2\_V3, che pur beneficiando della soglia (precisione da $\sim$.33 a $\sim$.459) resta la configurazione più problematica, con un numero di falsi positivi ancora elevato (483); ciò conferma che la sovrapposizione tra le classi V2 e V3 rappresenta il limite più significativo del modello, non completamente risolvibile con la sola modifica della soglia di confidenza.

\paragraph{Confronto con lo stato dell'arte}
All'interno della Tabella~\ref{tab:sota_confronto} vengono riportate le prestazioni raggiunte dal modello proposto (configurazione \emph{china modificato} con soglia 90\% dove mostrato) a fianco dei valori riportati in letteratura dai metodi classici di analisi del segnale e da modelli di deep learning. È importante premettere che si tratta di un confronto \emph{indicativo} e non di un benchmark diretto: i lavori citati utilizzano database, protocolli di simulazione dello scambio e definizioni di "successo" (per singolo battito, per finestra temporale o per intera registrazione) differenti tra loro e rispetto a quanto utilizzato in questa tesi; inoltre alcune metriche in letteratura sono aggregate su più tipologie di scambio contemporaneamente.

\begin{table}[htbp]
	\centering
	\caption{Confronto tra il modello proposto alcuni riferimenti di letteratura}
	\label{tab:sota_confronto}
	\rowcolors{2}{gray!15}{white}
	\scriptsize
	\begin{tabular}{p{3.3cm} p{2.6cm} p{2.6cm} p{2cm}}
		\toprule
		Riferimento & Tipo di metodo & Classe/i & Metrica principale \\
		\midrule
		Hedén et al. \cite{heden1995ann} & ANN (storica) & RA--LA & Sp$\approx$99.9\%, Se$\approx$98.7\% \\
		Kors \& van Herpen \cite{kors2001accurate} & Correlazione derivazioni & RA--LA/RA--LL & Se$\ge$93\%, Sp$\ge$99.5\% \\
		Kors \& van Herpen \cite{kors2001accurate} & Correlazione derivazioni & LA--LL & Se=17.9\% \\
		Jekova et al. \cite{jekova2016interlead} & Correlazione inter-lead & Precordiali (1--3 hop) & Se=99.52\%, Sp=99.62\% \\
		Jekova et al. \cite{jekova2016interlead} & Correlazione inter-lead & Periferiche & Se=92.10\%, Sp=96.05\% \\
		Rjoob et al. \cite{rjoob2021jmir} & Deep learning (BLSTM) & V1/V2 & Acc.=93.0\% \\
		Huang et al. \cite{huang2024lightweight} & Deep learning (lightweight) & arti+precordiali (12 casi) & Macro F1 93.42--99.61\% \\
		Alves et al. \cite{alves2025telecardiology} & Deep learning (CNN+RNN) & RA--LA, RA--LL & metriche $\geq$0.96 \\
		Alves et al. \cite{alves2025telecardiology} & Deep learning (CNN+RNN) & Precordiali & metriche $>$0.86 \\
		\midrule
		\textbf{Questo lavoro} & CNN & RA\_LL & F1=0.978 \\
		\textbf{Questo lavoro} & CNN & RA\_LA & F1=0.966 \\
		\textbf{Questo lavoro} & CNN & LA\_LL & F1=0.957 \\
		\textbf{Questo lavoro} & CNN & V1\_V2 & F1=0.943 \\
		\textbf{Questo lavoro} (soglia 90\%) & CNN & V3\_V4 & Precisione=0.833 \\
		\textbf{Questo lavoro} (soglia 90\%) & CNN & V2\_V3 & Precisione=0.459 \\
		\bottomrule
	\end{tabular}
\end{table}

\begin{table}[htbp]
	\centering
	\caption{Confronto tra il modello proposto (Se/Sp, seed 42) e alcuni riferimenti di letteratura}
	\label{tab:sota_confronto_sp_se}
	\rowcolors{2}{gray!15}{white}
	\scriptsize
	\begin{tabular}{p{3.2cm} p{2.4cm} p{2.2cm} p{1.4cm} p{1.4cm}}
		\toprule
		Riferimento & Tipo di metodo & Classe/i & Se & Sp \\
		\midrule
		Hedén et al. \cite{heden1995ann} & ANN (storica) & RA--LA & $\approx$98.7\% & $\approx$99.9\% \\
		Kors \& van Herpen \cite{kors2001accurate} & Correlazione derivazioni & RA--LA/RA--LL & $\geq$93\% (stimato) & $\geq$99.5\% \\
		Kors \& van Herpen \cite{kors2001accurate} & Correlazione derivazioni & LA--LL & 17.9\% & n.d. \\
		Jekova et al. \cite{jekova2016interlead} & Correlazione inter-lead & Precordiali (1--3 hop) & 99.52\% & 99.62\% \\
		Jekova et al. \cite{jekova2016interlead} & Correlazione inter-lead & Periferiche & 92.10\% & 96.05\% \\
		\midrule
		\textbf{Questo lavoro} & CNN & LA\_LL & 93.62\% & 99.98\% \\
		\textbf{Questo lavoro} & CNN & RA\_LA & 100.00\% & 99.93\% \\
		\textbf{Questo lavoro} & CNN & RA\_LL & 100.00\% & 99.95\% \\
		\textbf{Questo lavoro} & CNN & V1\_V2 & 99.29\% & 99.89\% \\
		\textbf{Questo lavoro} & CNN & V2\_V3 & 99.53\% & 97.97\% \\
		\textbf{Questo lavoro} & CNN & V3\_V4 & 98.82\% & 99.50\% \\
		\textbf{Questo lavoro} & CNN & V4\_V5 & 100.00\% & 98.86\% \\
		\textbf{Questo lavoro} & CNN & V5\_V6 & 99.29\% & 98.38\% \\
		\bottomrule
	\end{tabular}
\end{table}

\newpage
Sulle classi periferiche (RA\_LA, RA\_LL, LA\_LL), il modello proposto si colloca in linea con i risultati più recenti basati su deep learning (Alves et al.~\cite{alves2025telecardiology}) e supera nettamente il limite storico dei metodi classici sulla classe LA\_LL, dove Kors e van Herpen~\cite{kors2001accurate} riportavano una sensibilità di appena il 17.9\%: la rete convoluzionale qui proposta raggiunge invece un F1 di 0.957, con precisione e recall entrambi superiori al 92\%. Questo suggerisce che l'approccio basato su rappresentazioni apprese dal segnale grezzo, rispetto a regole di correlazione fissate a priori, sia particolarmente efficace proprio sui casi in cui i metodi classici mostravano più difficoltà.

Sulle classi precordiali adiacenti (in particolare V2\_V3), il confronto è invece meno favorevole: Jekova et al.~\cite{jekova2016interlead} riportano sensibilità e specificità superiori al 99\% per gli scambi precordiali, molto migliore rispetto alla precisione ottenuta dal modello proposto su questa classe (anche dopo la modifica della soglia di confidenza). Va tuttavia notato che il metodo di Jekova et al. sfrutta esplicitamente la struttura di correlazione nota tra derivazioni precordiali adiacenti, un'informazione strutturale che nel presente lavoro non viene fornita esplicitamente al modello ma dev'essere appresa implicitamente dai dati; questo potrebbe spiegare, almeno in parte, la netta differenza tra i risultati e suggerisce una possibili idea di miglioramento futuro (ad esempio incorporando vincoli o feature basate sulla correlazione inter-derivazione).

